%
% Documento: Conclusão
%

\chapter{Considerações finais}
\label{cap:conclusao}

Este trabalho substituiu a concepção inicial baseada em NFC por uma arquitetura RFID UHF mais compatível com a leitura de múltiplos bens durante uma varredura. Foram especificados o YPD-R200, a antena APCA8090, a etiqueta AD-239 M730 e o Arduino, e o firmware foi adaptado para a abstração \texttt{Stream}, comando nominal de potência e EPCs de tamanho variável.

A validação física planejada não foi realizada. A decisão de empregar simulação decorreu do tempo insuficiente para aquisição, montagem, ensaios piloto e repetições controladas e do custo de múltiplas antenas, etiquetas, instrumentação, estruturas de posicionamento e acesso a um ambiente controlado. Consequentemente, todas as taxas, RSSIs, rotas, passagens e janelas apresentadas são dados sintéticos. Eles não comprovam o desempenho do protótipo real.

No modelo, a taxa média simulada de leitura foi de 96,9\% a 0,5~m, 95,1\% a 1,0~m, 93,1\% a 1,5~m e 90,0\% a 2,0~m, caindo para 77,8\% a 2,5~m e 48,0\% a 3,0~m. Mesmo a 1,5~m, somente 10 das 15 configurações tag--altura atingiram individualmente o critério interno de 90\%, mostrando que a média da simulação pode ocultar condições desfavoráveis.

As taxas geradas para orientação foram de 89,1\%, 53,8\% e 10,0\% a $0^\circ$, $45^\circ$ e $90^\circ$. Essa queda foi produzida pela penalidade angular do modelo e não constitui medição da APCA8090. Para materiais, a simulação gerou 3,3\% sobre metal direto e 30,0\% com espaçador de 10~mm, em contraste com 80,0\% a 100,0\% nas superfícies não metálicas representadas.

No inventário virtual, a cobertura média foi de 91,0\%, com omissão de 9,0\%, mas somente três das dez janelas foram completas. Também foram geradas 298 duplicatas e duas leituras externas. Esses resultados exercitam as métricas e reforçam conceitualmente a necessidade de contar EPCs únicos; não demonstram confiabilidade operacional nem comprovam que o firmware atual realize essa consolidação.

A pergunta de pesquisa foi respondida somente no domínio do modelo: sob suas premissas, distância, orientação e material alteram fortemente as taxas projetadas, e a média de cobertura não basta para definir um inventário confiável. A hipótese apresentou coerência interna, mas não foi validada empiricamente. O trabalho deve ser entendido como arquitetura de protótipo acompanhada de análise exploratória e protocolo para futura experimentação.

Quanto à antena, uma alternativa omnidirecional pode ser mais adequada à varredura portátil quando a prioridade é ampliar a cobertura azimutal e reduzir o apontamento pelo operador. A simulação atual, entretanto, não comparou antenas e não permite declarar sua superioridade. Para reduzir a dependência da orientação das tags, a comparação mais direta continua sendo uma antena circular ou com diversidade de polarização. A escolha futura deve avaliar conjuntamente padrão de irradiação, polarização, ganho, leituras externas e ergonomia; para um portal fixo, o controle espacial de uma antena direcional pode ser preferível.

\section{Limitações e continuidade}

A principal limitação é o caráter integralmente sintético da base. O modelo é fenomenológico, não eletromagnético, e não foi calibrado com medições. Ele simplifica multipercurso, resposta de materiais, variação entre etiquetas, geometria, protocolo Gen2 e falhas da interface serial. Os percentuais dependem das regras adotadas e não devem ser generalizados para o YPD-R200, a APCA8090 ou outras instalações.

Embora a planilha registre sementes e preserve todas as 6.330 tentativas controladas e os 400 registros de inventário, o código original do gerador, seus coeficientes e a função de conversão para probabilidade não estão no repositório. Assim, os agregados são auditáveis a partir da base, mas a geração não é integralmente reproduzível do zero. Uma evolução prioritária é versionar o gerador completo ou reconstruir um modelo documentado e calibrá-lo com uma coleta piloto.

O firmware também permanece incompleto para inventário real: armazena apenas o último EPC, não mantém uma tabela de EPCs únicos por janela, não grava RSSI bruto, timestamp e condição experimental e não confirma região, canal, FHSS ou potência aplicada. Esses requisitos devem ser implementados antes de qualquer comparação entre dados físicos e simulados.

Como continuidade, recomenda-se realizar inicialmente uma campanha reduzida e financeiramente viável, preservando log serial integral, fotografias, geometria, versão do firmware e configuração do leitor. Em seguida, deve-se comparar a antena linear de referência com uma alternativa circular ou de diversidade de polarização e uma alternativa omnidirecional, usando as mesmas tags, cabos, potência e ambiente. Para bens metálicos, devem ser incluídas etiquetas próprias para metal. Somente após essa calibração será possível revisar o modelo, quantificar sua aderência e concluir sobre o desempenho operacional do protótipo.
